\documentclass[10pt,a4paper]{amsart}
\usepackage[margin=19mm]{geometry}
\usepackage{amsmath,amssymb,mathtools}
\usepackage[scr=rsfs,cal=euler]{mathalfa}
\usepackage{xfrac}
\usepackage[unicode,hidelinks,bookmarks=false]{hyperref}
\newcommand{\ie}{i.e.\ }
\newcommand{\cf}{cf.\ }
\newcommand{\resp}{respectively\ }
\newcommand{\Subsection}{\subsection}
\providecommand{\nobreakdash}{\nobreak\hbox{-}}
\setlength{\emergencystretch}{2em}
\multlinegap0pt
\usepackage{bm}
\usepackage{bbm}
\usepackage{dsfont}
\usepackage{xspace}
\usepackage{cancel}
\usepackage{mathtools}
\usepackage{amsthm}
\makeatletter
\@ifundefined{theorem}{\newtheorem{theorem}{Theorem}}{}
\@ifundefined{lemma}{\newtheorem{lemma}[theorem]{Lemma}}{}
\makeatother
\title[The Sandpile Group of a Cone Over a Bi-Coconut Tree]{The Sandpile Group of a Cone Over a Bi-Coconut Tree}
\author{Dorian Smith}
\date{Source: arXiv:2602.18621v1 (CC BY 4.0)}
\begin{document}
\maketitle

\noindent\emph{Source and licence.} The Sandpile Group of a Cone Over a Bi-Coconut Tree. Version arXiv:2602.18621v1 (CC BY 4.0), distributed under the Creative Commons Attribution 4.0 International licence, as recorded in the arXiv licence metadata for this exact version and shown on the abstract page https://arxiv.org/abs/2602.18621v1. Full rights notice: licenses/arxiv-2602.18621-CC-BY-4.0.txt.\par\smallskip

\noindent\textit{Editorial scope.} Counted unit: the Lemma whose source label is \texttt{a-is-2x} in the article (source0.tex lines 1604--1608, source label \texttt{a-is-2x}), delivered together with the article's own complete original proof of it (source0.tex lines 1610--1680, one \texttt{proof} environment of 71 lines). No mathematics is written, repaired or re-derived: statement and proof are the article's own text line for line. Required context retained uncounted: formula-for-x (lines 1582--1586); formula-for-a-repeated (lines 1596--1599). Every definition, display and label of those lines is retained unchanged. Omitted from this package and not redistributed: 1--1581, 1587--1595, 1600--1603, 1609--1609, 1681--2337. Nothing inside a retained excerpt is omitted or altered: every retained body is delivered exactly as the article's lines are written, so no omission receipt of any kind accompanies this package. Citations: the retained excerpts carry no citation command at all, so no bibliography entry of the article is transcribed and no reference is redistributed. Reference adaptation: none was needed and none was made; every reference inside the delivered unit resolves inside it, so no unresolved reference or citation is produced. Printed slips kept literal: no printed word, symbol, hypothesis or number of the counted statement or of its proof was altered. Layout and class: the article's own class is \texttt{amsart} with packages including inputenc, float, multicol, fullpage, tikz-cd, graphicx, enumitem, amsmath, amssymb, amsthm, amsfonts, verbatim, pdfpages, wrapfig; the wrapper is the lane's pinned \texttt{amsart} editorial wrapper at 10pt on A4, which restates the theorem-like environments the retained text needs and adds the article's own macro definitions that the retained text needs (none), with extra wrapper packages (bm, bbm, dsfont, xspace, cancel, mathtools). The delivered numbering is the unit's own. Assets: no retained span contains a figure environment, a graphics-inclusion command, a PDF-page inclusion, a table or a TikZ picture; no image file is copied into this package, the package declares no asset dependency and no image file is redistributed with it. Licence: version arXiv:2602.18621v1 of this article is distributed under the Creative Commons Attribution 4.0 International licence, as recorded in the arXiv licence metadata for this exact version and shown on the abstract page https://arxiv.org/abs/2602.18621v1. A full rights notice is delivered at licenses/arxiv-2602.18621-CC-BY-4.0.txt. Characters: every retained excerpt is pure ASCII, so the delivered engine, XeLaTeX with its default Latin Modern text font, reports no missing character.
\begin{align} 
\label{formula-for-x}
 \notag   x(p)&:= 2(df+e) - s_2d-s_1(gf+h-\frac{1}{2}s_2g),\\
    &= (2s_1+2s_2+\displaystyle\frac{1}{2}s_1s_2+8)F_{2p-2}-(s_1+s_2+8)F_{2p-4}+2F_{2p-6}.
\end{align}

\begin{equation}
\label{formula-for-a-repeated}
a:=a(p):=2^{1-s_1}(2\cdot b_{2p-3} +s_2b_{2p-4}).
\end{equation}

\begin{lemma}
\label{a-is-2x}
Let $p \geq 2.$ The two quantities $x(p)$ from \eqref{formula-for-x} and $a(p)$ from \eqref{formula-for-a-repeated} are related by
$x(p) = \displaystyle\frac{a(p)}{2}.$
\end{lemma} 

\begin{proof}
Moving forward we let 
\begin{align}
    A & = 2s_1+2s_2+\displaystyle\frac{1}{2}s_1s_2+8.\notag\\
    B &= -(s_1+s_2+8)\notag.\\
    C &= 2^{1-s_1}.\notag\\
D &= 2^{-s_1}s_2.\notag
\end{align}
Observe that if we let
\begin{align}
    x'(p)&= AF_{2p-4}+BF_{2p-6}+2F_{2p-8}\notag\\
    l(p)&=Cb_{2p-5}+Db_{2p-6}\notag\\
y'(p)&= AF_{2p-3}+BF_{2p-5}+2F_{2p-7}\notag\\
m(p)&=Cb_{2p-4}+Db_{2p-5}.\notag
\end{align}
The formulas \eqref{formula-for-x}, \eqref{formula-for-a-repeated} show that
the two sides of the lemma can both be rewritten as sums of two terms 
\begin{align}   
x(p)&=x'(p)+y'(p)\\
\frac{a(p)}{2}&=l(p)+m(p)
\end{align} 
hence, the lemma asserts the following:
\begin{equation}
   x'(p)+y'(p) =l(p)+m(p) \label{x'_y'} \hspace{1cm}
%\text{iff}\\  x(p) &= \displaystyle\frac{a(p)}{2}\notag.
\end{equation}
Thus it suffices to prove simultaneously by induction on $p$ that

\begin{align}    
x'(p) &= l(p) \label{x and l} \hspace{1cm} \text{and}\\
y'(p) &= m(p) \label{y and m}  
\end{align}
Base case: We check when $p = 2, p=3, p=4$ we see that 
\begin{align}  
x'(2) &= l(2) = s_1 +s_2 + 2.\notag\\
y'(2) &= m(2)   = s_1+ s_2+\displaystyle\frac{1}{2}s_1s_2+4\notag\\
x'(3) &= l(3) = 2s_1 +2s_2 +\displaystyle\frac{1}{2}s_1s_2+6\notag\\
y'(3) &= m(3)   = 3s_1+ 3s_2+s_1s_2+10\notag\\
x'(4) &= l(4) = 16+5s_1 +5s_2 + \displaystyle\frac{3}{2}s_1s_2.\notag\\
y'(4) &= m(4)   = 26+8s_1+8s_2+\displaystyle\frac{5}{2}s_1s_2\notag\\
\end{align}
Now assume that the statements  \eqref{x and l} and \eqref{y and m}
are true for $p = k, k-1, k-2.$ That is, 

\begin{align}
 AF_{2k-4}+BF_{2k-6}+2F_{2k-8}
   &=Cb_{2k-5}+Db_{2k-6} \label{ABCD_1}
   \end{align} and 
 \begin{align} 
    AF_{2k-3}+BF_{2k-5}+2F_{2k-7}&=Cb_{2k-4}+Db_{2k-5}.\label{ABCD_2}
\end{align} 

Then the sum of equations \eqref{ABCD_1}  and \eqref{ABCD_2}  is true and equal to  \eqref{x'_y'} evaluated at $k+1:$  
\begin{align}
 AF_{2k-2}+BF_{2k-4}+2F_{2k-6}
   &=Cb_{2k-3}+Db_{2k-4}.\notag
   \end{align}
Next our inductive step yields the following:

\begin{align}
 Cb_{2k-2}+Db_{2k-2}&=Cb_{2k-3}+Cb_{2k-4}+Db_{2k-4}+Db_{2k-5}\notag\\
 &=Cb_{2k-4}+Cb_{2k-5}+Cb_{2k-4}+Db_{2k-5}+Db_{2k-6}+Db_{2k-5}\notag\\
 &=AF_{2k-4}+BF_{2k-6}+2F_{2k-8}\notag\\
 & \qquad +AF_{2k-3}+BF_{2k-5}+2F_{2k-7}+AF_{2k-3}+BF_{2k-5}+2F_{2k-7}\notag\\
 &=A(F_{2k-4} + F_{2k-3}+F_{2k-3})\notag\\
 &\notag 
 \qquad +B(F_{2k-5}+ F_{2k-6} +F_{2k-5})+2(F_{2k-7}+F_{2k-8}+F_{2k-7})\\
 &=AF_{2k-1}+BF_{2k-3}+2F_{2k-5} \notag
   \end{align}
   which is equal to \eqref{ABCD_1} evaluated at $k+1$.
\end{proof}

\end{document}
